\documentstyle[%  


righttag%  


,11pt%  


,amssymb%  


,russian%               remove in Eng.  


,izvru%                 Set izven% in Eng.  


]{amsart}  


  


% 


\newcommand{\la}{\langle}  


\newcommand{\ra}{\rangle}  


\newcommand{\supp}{\operatorname{supp}}  


\newcommand{\re}{\operatorname{Re}}  


\newcommand{\sgn}{\operatorname{sgn}}  


\newcommand{\tts}[1]{}  


  


\theoremstyle{plain}  


\theorembodyfont{\it}  


\newtheorem{thms}{\hskip\parindent Теорема}[section]  


\newtheorem{lems}{\hskip\parindent Лемма}[section]  


\theoremstyle{definition}  


\theorembodyfont{\rm}  


\newtheorem{notes}{\hskip\parindent Замечание}[section]  


\numberwithin{equation}{section}  


\setcounter{section}{2} 


\setcounter{equation}{19} 


\setcounter{notes}{1} 


 


%\hoffset=-15mm%                  For Eng.  


\setcounter{page}{30}  


\god{1999}  


\nomer{7~(446)} %                        Remove in Eng.  


%\nomer{Vol.\,43, No.\,7}%               Set in Eng.  


%\str{\pageref{begin}--\pageref{end}}%  Set in Eng.  


\udk{517.9}  


 


\author{\it В.С.\,Мокейчев, \fbox{А.В.\,Мокейчев}}  


% NB:   ^^ remove in Eng.  


\title{Новый подход к теории линейных задач\\  


для систем дифференциальных уравнений\\  


в частных производных, II}  


  


\date{}  


% \pagestyle{myheadings}%         Set in Eng.  


\pagestyle{plain} %              Remove in Eng.  


\begin{document}  


% \thispagestyle{plain}%          Set in Eng.  


\maketitle  


\label{begin}  


 


Мы сохраняем обозначения, понятия, соглашения и нумерацию, 


использованные 


в [1]. 


 


Пусть $u$ --- $\varphi$-решение линейной задачи (7), (8) т.\,е. задачи 


$$ 


(P(x,\partial/\partial x)-\lambda E)u=f(x)\in H,\quad 


x\in\Omega\subset\Bbb R^n;\quad l u=0, 


$$ 


и $u_k$ --- вектор-столбец с координатами $u_{k,j}$, $j=1,\dotsc,m$, 


являющимися 


коэффициентами Фурье с номерами $(k,j)$ по последовательности 


$\varphi=\{e(j)y_{(k)},\ k\in\cal N,\ j=1,\dotsc,m\}$. 


Если функции $y_{(k)}$ имеют производные порядка  


$\beta$, и ряд $\sum u_k y_{(k)}^{(\beta)}$ сходится по норме в некотором 


гильбертовом 


пространстве $H_{(\beta)}$, то его сумму $u^{(\beta)}$ следует назвать 


производной 


порядка $\beta$ от $\varphi$-рас\-пре\-де\-ле\-ния~$u$.\; 


$\varphi$-распределение $u$ 


может быть дифференцируемым в $H_{(\beta)}$, но недифференцируемым в 


другом  


гильбертовом пространстве. Если $\varphi$-распределение $u$ имеет в 


$H_{(\beta)}$ 


производную порядка $\beta$ и $\alpha<\beta$, то даже в случае 


существования $u^{(\alpha)}$ 


включение $u^{(\alpha)}\in H_{(\beta)}$ может не выполниться. Простой 


пример: $\varphi$ 


совпадает c (1.4), $\beta=(1,1)$, $\alpha=(1,0)$, 


$H_{(\beta)}=L^2_m(b-a)$ и 


$$ 


u=\sum_{k_1}u_{k_1,0}\exp\bigg(2\pi i k_1 


\frac{x_1-a_1}{b_1-a_1}\bigg)+\sum_{k_2}u_{0,k_2}\exp 


\bigg(2\pi i k_2\frac{x_2-a_2}{b_2-a_2}\bigg). 


$$ 


 


{\bf Теорема 2.3 $($o глобальной регулярности$)$.} {\it Пусть выполняются  


оценки $(2.1)$. Тогда следующие утверждения равносильны$:$ 


\begin{itemize} 


\item[1.] для каждой $f(x)$ в случае существования $\varphi$-решения 


задачи  


$(7)$, $(8)$ оно имеет в $H_{(\beta)}$ производную порядка $\beta;$ 


\item[2.] при всех $\psi\in L_\varphi$ выполняются оценки  


\begin{equation} 


c_4\|\psi^{(\beta)}\|_\beta-c(N)\|\psi(N,x)\|\le 


\|Q\psi\|, 


\end{equation} 


в которых $c_4>0$, $N$ и $c(N)$ не зависят от $\psi$ и 


$\|\cdot\|_\beta$ --- символ нормы 


в $H_{(\beta)}$.


\end{itemize} 


} 


 


\begin{pf} 


Напомним, что оценки (2.1) имеют вид 


$$ 


c_1\|Q\psi\|-c(N)\|\psi(N,x)\|\le\|P_\lambda\psi\|\le c_2 


\|Q\psi\|+c(N)\|\psi(N,x)\|, 


$$ 


и $Q$ --- $\varphi$-стационарный, взаимно однозначный оператор, 


определяемый 


равенствами $Qw=\sum{}q(k)w_k y_{(k)}$, в которых $q(k)$ --- обратимые 


матрицы, 


не зависящие от $x$. Пусть имеет место утверждение 2, и $u$ --- 


$\varphi$-решение задачи (7), (8). В силу (2.1) в $H$ по норме сходится 


ряд 


$$ 


\sum Q[u_k y_{(k)}]=\sum q(k)u_k y_{(k)}. 


$$ 


Отсюда и из оценок, используемых в утверждении 2, следует сходимость 


по норме в $H_{(\beta)}$ ряда $\sum u_k y^{(\beta)}_{(k)}$. А это 


означает, что $u$ имеет в 


$H_{(\beta)}$ производную $u^{(\beta)}$. Импликация $2\Rightarrow 1$ 


доказана.


 


Докажем обратную импликацию. 


В силу утверждения 1 и оценок (2.1) ряд $\sum u_k y^{(\beta)}_{(k)}$ 


сходится 


по норме в $H_{(\beta)}$, если в $H$ по норме сходится ряд 


$\sum q(k)u_k y_{(k)}$. 


Отсюда и из ортонормированности $\varphi$ следует сходимость по норме в 


$H_{(\beta)}$ ряда $\sum(q(k))^{-1}v_k y^{(\beta)}_{(k)}$ при условии 


$\sum|v_k|^2<+\infty$. Сказанное позволяет 


определить оператор $T$ равенствами 


$$ 


T w=\sum(q(k))^{-1}w_k y^{(\beta)}_{(k)}\quad 


\forall w\in H. 


$$ 


Напомним, что $w_k$ --- вектор с координатами $w_{k,j}$, 


$j=1,\dotsc,m$, и $w_{k,j}$ --- 


коэффициент Фурье с номером $(k,j)$\; $\varphi$-распределения $w$.


Значит, $T$ определен на всем $H$. Докажем его замкнутость. 


Пусть $g\in H_{(\beta)}$. Тогда 


$$ 


\la T w,g\ra_\beta=\sum w_k\bullet(((q(k))^*)^{-1}g_k). 


$$ 


Здесь и ниже $\la\ ,\ \ra_\beta$ --- символ скалярного произведения в 


$H_{(\beta)}$, $\bullet$ --- 


символ скалярного произведения в $\Bbb C^m$. Так как последний ряд 


сходится 


при любых числовых векторах $w_k$, удовлетворяющих условию 


$\sum|w_k|^2<+\infty$,  


то в силу известной резонансной теоремы Э.\,Ландау 


$$ 


\gamma_0=\sum|((q(k))^*)^{-1}g_k|^2<+\infty. 


$$ 


Предположим, что $\|T w(\nu)-v\|_\beta+\|w(\nu)-h\|\to 0$ при 


$\nu\to+\infty$. Тогда 


$$ 


|\la T(w(\nu)-h),g\ra_\beta|\le\sqrt{\gamma_0}\big(\sum 


|(w(\nu))_k-h_k|^2\big)^{1/2}. 


$$ 


Поэтому 


$|\la T(w(\nu)-h),g\ra_\beta|\le\sqrt{\gamma_0}\|w(\nu)-h\|\to 0$ 


при $\nu\to+\infty$. Другими словами, 


доказано, что $T(w(\nu)-h)\to 0$ слабо. Однако $T w(\nu)\to v$ по 


норме, поэтому  


$T h=v$. Итак, доказано, что $T$ --- замкнутый оператор, определенный 


на всем $H$. В силу теоремы Банаха о замкнутых операторах $T$ 


непрерывен, т.\,е. 


$\big\|\sum(q(k))^{-1}w_k y^{(\beta)}_{(k)}\big\|_\beta\le\|T\|\,\|w\|$. 


Отсюда легко следуют оценки (2.20). 


\end{pf} 


 


{\bf Теорема 2.4 $($об индивидуальной регулярности$)$.} {\it Пусть 


$u$ --- 


$\varphi$-решение задачи $(7)$, $(8);$ выполняются оценки $(2.1);$ 


операторы 


$c_\alpha(x)$, $\alpha\in\Phi$, и функция $f(x)$ настолько гладки, что 


при некотором $\varphi$-стационарном 


операторе $\wt Q$, взаимно однозначно отображающем $\cal D_{\Tilde 


Q}\subset H$ на $H$, 


справедливо включение $u\in\cal D_{\Tilde Q}$ и имеет $\varphi$-решение 


уравнение 


\begin{equation} 


[\wt Q(P(x,\partial/\partial x)-\wt\lambda E)]v=\wt Q f(x)- 


\wt\lambda\wt Q u, 


\end{equation} 


в котором $\wt\lambda$ --- несобственное значение оператора 


$P_{\Tilde\lambda};$ при некоторой 


постоянной $d_1<c_1$ $($см. $(2.1))$ и всех $\psi\in L_\varphi$ 


\begin{equation} 


\|(\wt Q P_{\Tilde\lambda}-P_{\Tilde\lambda}\wt Q)\psi\|\le 


d_1\|P_{\Tilde\lambda}\wt Q\psi\|+c(N)\|\psi(N,x)\|. 


\end{equation} 


Тогда 


\begin{equation} 


\sum|q(k)\wt q(k)u_k|^2<+\infty, 


\end{equation} 


где $q(k)$, $\wt q(k)$ --- матрицы, определяемые 


$\varphi$-стационарностью 


соответственно $Q$, $\wt Q$. 


}


 


\begin{pf} 


$\varphi$-решение уравнения (2.21) обозначим через $v$.  


В силу (2.1), (2.22) его коэффициенты Фурье $v_{k,j}$ удовлетворяют 


(2.23). 


По определению $\varphi$-решения имеем 


$$ 


\sum\wt Q(P(x,\partial/\partial x)-\wt\lambda E)[v_k y_{(k)}]= 


\wt Q f(x)-\wt\lambda\wt Q u, 


$$ 


и ряд сходится в $H$ по норме. Однако оператор $\wt Q^{\,-1}$ 


непрерывен, поэтому 


$$ 


\sum(P(x,\partial/\partial x)-\wt\lambda E)[v_k y_{(k)}]= 


f(x)-\wt\lambda u, 


$$ 


и ряд сходится в $H$ по норме. Следовательно, $v$ --- $\varphi$-решение 


уравнения $(P(x,\partial/\partial x)-\wt\lambda E)z=f(x)-\wt\lambda u$, 


причем $u$ --- также его $\varphi$-решение. 


Так как $\wt\lambda$ --- несобственное значение оператора 


$P_{\Tilde\lambda}$, то $u=v$. 


\end{pf} 


 


Применим полученные результаты к математической модели (1), (2), 


введя в нее спектральный параметр $\lambda$. Итак, 


\begin{equation} 


u_{tt}-c^2u_{xx}-\lambda u=f(t,x),\quad 


H=L^2_1(\Bbb R\times(0,\pi)). 


\end{equation} 


 


Спектр $\varphi$-задачи для (2.24) в случае 


$$ 


\varphi=\big\{\tfrac{1}{\pi}\exp(i k_1t)\sin(k_2x), 


\ k_1=0,\pm1,\dotsc,\ k_2=1,2,\dotsc\big\} 


$$ 


совпадает с множеством 


$\sigma=\{-k^2_1+c^2k^2_2,\ k_1=0,1,\dotsc,\ k_2=1,2,\dotsc\}$. Если 


$\lambda\notin\sigma$, то $\varphi$-задача однозначно разрешима, и ее 


решение имеет вид 


$$ 


\sum(f_{k_1,k_2}/(-k^2_1+c^2k^2_2))\exp 


(i k_1t)\sin(k_2x). 


$$ 


Оно будет принадлежать $H$ при любой $f(x)\in H$ тогда и только тогда, 


когда 


$|k^2_1-c^2k^2_2|\ge\delta_1>0$ $\forall k_1,k_2$. 


 


\begin{notes} 


Выше предполагалось, что оператор $Q$ 


взаимно однозначен.  


Без этого предположения теоремы 2.1--2.4 справедливы только тогда,  


когда множество всех линейно независимых решений уравнения $Qv=0$ 


конечно. Если последнее множество бесконечно, то каждое 


$\lambda\in\Bbb C$ является собственным для $P_\lambda$, при этом для 


сохранения теоремы 2.3 нужно в 


утверждении 2 добавить требование: множество 


$\cal N_\beta= 


\{k:\det(q(k))=0\}\cap\{r:\|e(1)y^{(\beta)}_{(r)}\|\ne0\}$ 


конечно. 


\end{notes} 


 


\section{Связь $\varphi$-решений с $\wt\varphi$-решениями} 


 


Предположим, что для задачи (7), (8) мы можем вычислить 


$\varphi$-решение 


$u$ и $\wt\varphi$-решение $v$. Возникает вопрос: какова связь объектов 


$u$ и $v$?  


Чтобы ответить на поставленный вопрос, обозначим через $P_\lambda$, 


$\wt P_\lambda$ операторы, 


порожденные соответственно $\varphi$- и $\wt\varphi$-задачами для (7), 


через $Q$, 


$\wt Q$ соответственно $\varphi$- и $\wt\varphi$-стационарные операторы. 


При этом будем  


предполагать, что последние два оператора взаимно однозначны и 


$\cal D_{P_\lambda}=\cal D_Q$, $\cal D_{\Tilde P_\lambda}=\cal 


D_{\Tilde Q}$. 


Напомним, что $\cal D_B$ --- область определения оператора $B$. 


В общем случае $\cal D_B\subset\Bbb D\,'_\varphi$ либо $\cal 


D_B\subset\Bbb D\,'_{\Tilde\varphi}$. Тот факт, что из оценок (2.1) 


следует равенство $\cal D_{P_\lambda}=\cal D_Q$, почти очевиден. Чтобы 


убедиться в обратном, 


достаточно повторить рассуждения, аналогичные тем, которые были 


использованы в доказательстве импликации $1\Rightarrow 2$ 


в теореме 2.3. 


 


\begin{thms} 


Предположим, что $\cal D_{P_\lambda}=\cal D_Q$, $\cal D_{\Tilde 


P_\lambda}=\cal D_{\Tilde Q}$, последовательности $\varphi$ и 


$\wt\varphi$ 


являются ортобазисами в гильбертовом пространстве $H_{(0)}$, и  


при каждом $p\in\cal N$ выполняются равенства 


\begin{equation} 


P_\lambda\varphi_{(p)}=\sum a_{p,k}\wt P_\lambda\wt\varphi_{(k)}, 


\end{equation} 


причем ряды сходятся в $H$ по норме. Тогда $\varphi$-решение $u$ 


уравнения $(7)$ 


продолжается до $\wt\varphi$-решения того же уравнения$;$ если, кроме 


того, 


$\cal D_Q\subset H_{(0)}$, $\cal D_{\Tilde Q}\subset H_{(0)}$, то $u$ и 


его продолжение $\wt u$ совпадают как элементы из $H_{(0)}$. 


\end{thms} 


 


\begin{pf} 


В частном случае для конкретных $P$ и $\varphi$ теорема доказана 


в ([2], c.\,196--200). Докажем ее в общем случае. В силу (3.1) область 


значений $\cal J_1$ оператора $P_\lambda$ содержится в области значений 


$\cal J_2$ 


оператора $\wt P_\lambda$. Так как выполняются оценки (2.1) и их 


аналоги, получающиеся 


заменой в (2.1) $P_\lambda$ на $\wt P_\lambda$, $Q$ на $\wt Q$ и 


$L_\varphi$ на $L_{\Tilde\varphi}$, то в силу теоремы 2.1\; 


$\cal J_1$ и $\cal J_2$ --- подпространства в $H$, причем в силу взаимной 


однозначности $\wt Q$ последовательность 


\begin{equation} 


\{\wt P_\lambda(\wt Q)^{-1}\wt\varphi_{(p)}, 


\ p\in\wt{\cal N}\},\ \;\wt{\cal N}\subset\cal N,


\end{equation} 


является базисом Рисса в $\cal J_2$. Справедливость этого факта была 


установлена 


в процессе доказательства теоремы 2.2. При этом отмечалось,  


что в множестве $\cal N\setminus\wt{\cal N}$ 


число элементов конечно. Так 


как (3.2) --- базис 


Рисса в $\cal J_2$ и $\wt P_\lambda\wt\varphi_{(k)}\in\cal J_2$, то 


\begin{equation} 


\wt P_\lambda\wt\varphi_{(k)}=\sideset{}{'}\sum_p 


\beta_{k,p}\wt P_\lambda\wt\varphi_{(p)}\quad 


\forall k\in\cal N. 


\end{equation} 


Штрих указывает на то, что суммирование производится по всем 


$p\in\wt{\cal N}$. 


Из последних рассуждений следует, что наряду с (3.1) имеем


$$ 


P_\lambda\varphi_{(k)}=\sideset{}{'}\sum\wt a_{k,p}\wt P_\lambda 


\wt\varphi_{(p)}\quad \forall k\in\cal N. 


$$ 


Таким образом, заменяя, если в этом есть необходимость, $\wt\varphi$ на 


последовательность $\{\wt\varphi_{(p)},\ p\in\wt{\cal N}\}$, добьемся 


того, что $\lambda$ станет несобственным 


значением оператора $\wt P_\lambda$. Поэтому, не нарушая общности, 


можно считать, 


что $\lambda$ --- несобственное значение оператора $\wt P_\lambda$, 


т.\,е. $\wt{\cal N}=\cal N$. 


 


В силу (3.1), (3.2) имеем 


\begin{equation} 


P_\lambda u=\sum P_\lambda[u_p\varphi_{(p)}]=\sum_p 


\Big\{\sum_q a_{q,p}\wt P_\lambda\wt\varphi_{(q)}\Big\}u_p. 


\end{equation} 


Так как $\cal J_1\subset \cal J_2$, то существует такое 


$\wt\varphi$-распределение $v$, что 


\begin{equation} 


\wt P_\lambda v=P_\lambda u=\sum v_q\wt P_\lambda\wt\varphi_{(q)}. 


\end{equation} 


Однако (3.2) --- базис Рисса в $\cal J_2$. Поэтому равенства (3.4), (3.5) 


выполнятся тогда и только тогда, когда 


\begin{equation} 


v_q=\sum_p a_{p,q}u_p\quad \forall q\in\cal N. 


\end{equation} 


Докажем, что $\varphi$-решение $u$ можно продолжить до 


$\wt\varphi$-решения $\wt u$. 


Обозначим через $L$ линейную оболочку множества 


$\varphi\cup\wt\varphi$. В силу (3.6) 


и конечности множества $\cal N\setminus\wt{\cal N}$ для каждой 


$\psi\in{}L$ ряд $\sum u_p\la\psi,\varphi_{(p)}\ra$ сходится. 


Поэтому можно ввести функционал $\wt u$, задаваемый на $L$ с 


помощью равенств $\wt u(\varphi)=\sum u_p\la\psi,\varphi_{(p)}\ra$. 


Очевидно, выписанный 


функционал является как $\varphi$-распределением, так и 


$\wt\varphi$-распределением, 


причем $\wt u=u$ на $L_\varphi$, и в силу (3.6) $\wt u=v$ на 


$L_{\Tilde\varphi}$. 


Следовательно, 


функционал $\wt u$ является как $\varphi$-решением, так и 


$\wt\varphi$-решением уравнения (7).  


Первое утверждение теоремы доказано. Чтобы доказать второе, заметим 


$\wt u=u$ на $L_\varphi$, причем $\wt u\in H_{(0)}$ и $u\in H_{(0)}$. 


Однако $\varphi$ --- базис в $H_{(0)}$, 


поэтому $\|\wt u-u\|_0=0$. 


\end{pf} 


 


\section{Корректно разрешимая задача для  


дифференциальных квазиполиномов\protect\\


с постоянными коэффициентами} 


 


Пусть $P(z)$ --- символ $(m,m)$ матричного дифференциального 


квазиполинома $P(\cal D)$ 


с постоянными коэффициентами, т.\,е. при всех $z\in\Bbb C^m$, 


$x\in\Bbb R^n$ 


$$ 


P(\cal D)\exp(z\bullet x)=\bigg[\,\sum_{\alpha\in\Phi}\,\sum^M_{j=1} 


a_{\alpha,j}(z^\alpha)\exp(z\bullet\tau_{\alpha,j})\bigg] 


\exp(z\bullet x)\equiv P(z)\exp(z\bullet x), 


$$ 


где множество мультииндексов $\Phi$, числовые векторы $\tau_{\alpha,j}$ 


и числовые  


матрицы $a_{\alpha,j}$ не зависят от $z$, $x$, причем $\Phi$ конечно. 


 


Обозначим $\Vec\mu=\mu\cdot t$, где $\mu\in\Bbb C$ и $t\in\Bbb C^n$, 


$\sigma_k=\{\mu:\det(P(\Vec\mu+ik))=0\}$, 


$\sigma=\bigcup\limits_{k\in\Bbb Z^n}\sigma_k$. 


Очевидно, скалярная функция $\det(P(\Vec\mu+ik))$ при каждом $k$ 


аналитически зависит от $\mu$. Поэтому в случае 


$\det(P(\Vec\mu+ik))\not\equiv0$ 


множество $\sigma_k$ счетно. Если каждое множество $\sigma_k$ счетно, 


то и множество $\sigma$ 


счетно. При этом пустое множество считаем счетным. 


 


\begin{thms} 


Пусть при некотором $t$ и каждом $k$ имеет место 


$\det(P(\Vec\mu+ik))\not\equiv0$. Если $\mu\notin\sigma$, то в случае 


$$ 


\varphi=\{e(j)\exp((\Vec\mu+ik)\bullet x), 


\ k\in\Bbb Z^n,\ j=1,\dotsc,m\},\quad \Omega=(0,2\pi)^n 


$$ 


$\varphi$-задача для уравнения $P(\cal D)u=f\in L^2_m(\Omega)$ 


корректно разрешима.  


\end{thms} 


 


\begin{pf} 


Но определению $\varphi$-решения имеем 


$$ 


P(\cal D)u=\sum_{k\in\Bbb Z^n}P(\Vec\mu+ik)u_k\exp 


((\Vec\mu+ik)\bullet x)=f. 


$$ 


Так как последовательность $\varphi$ является базисом в 


$L^2_m(\Omega)$, то $f=\sum\limits_{k\in\Bbb 


Z^n}f_k\exp((\Vec\mu+ik)\bullet x)$. 


Поэтому $u_k=(P(\Vec\mu+ik))^{-1}f_k$. Следовательно, 


$\varphi$-задача однозначно разрешима. Векторы $u_k$ при каждом $k$ 


непрерывно 


зависят от $f_k$, причем последние непрерывно зависят от $f$. 


Отсюда следует непрерывная зависимость $\varphi$-решения $u$ от $f$. 


\end{pf} 


 


\section{Приложения} 


 


Что касается математической модели (1), (2), 


то о ней подробно говорилось в конце второго раздела. Всюду ниже,  


если особо не оговорено, $\|\cdot\|$ --- символ нормы в $H$, 


порожденной  


скалярным произведением $\la\ ,\ \ra$. 


 


1. {\it $\varphi$-задача в $L^2_m(\Bbb R^n)$}. Ниже предполагается, 


что $H=L^2_m(\Bbb R^n)$, 


$\varphi=\Big\{e(j)\prod\limits^n_{\nu=1}y_{(k_\nu,\nu)}(x_\nu), 


\ k_\nu\in\Bbb Z_+,\ j=1,\dotsc,m\Big\}$, для 


каждого 


$\nu$ последовательность $\{y_{(k_\nu,\nu)}(x_\nu),\ k_\nu\in \Bbb Z_+\}$ 


является ортонормированным 


базисом в $L^2_1(\Bbb R)$, причем 


\begin{equation} 


\begin{array}{c} 


y_{(r,\nu)}(t)\in W^{+\infty,2}_1(\Bbb R),\quad 


y^{(s)}_{(r,\nu)}(t)\to 0,\quad s=0,1,\dotsc,\ \; 


\text{при}\ \; |t|\to+\infty,\\ 


\displaystyle\bigg(\bigg(\frac{d}{dt}\bigg)^2+h_\nu(t)-% 


\lambda_{r,\nu}\bigg)y_{(r,\nu)}(t)=0\quad 


\forall t\in\Bbb R, 


\end{array} 


\end{equation} 


функции $h_\nu(t)\in C^{+\infty}_1(\Bbb R)$ удовлетворяют условиям 


$\re h_\nu(t)\le\alpha_0$, в которых $\alpha_0$ 


не зависит от $t$, $\lambda_{r,\nu}\ne 0$ $\forall r,\nu$. 


 


Простейшим примером функций $y_{(r,\nu)}(t)$, удовлетворяющих 


отмеченным выше 


условиям, являются функции из ортонормированной последовательности 


Чебышева--Эрмита [3]. 


 


Рассмотрим только те уравнения (7), которые представимы в виде 


\begin{multline} 


P_1(x,\partial/\partial x)u\equiv\sum_{\alpha\in\Phi(2)}\bigg[ 


a_\alpha(x)((\partial/\partial x)^2+h)^\alpha+ \\


%


+\sum_{\beta<\alpha}b_{\alpha,\beta}(x) 


(\partial/\partial x)^{\gamma(\alpha,\beta)} 


((\partial/\partial x)^2+h)^\beta\bigg]u-\lambda u=f(x), 


\quad x\in\Bbb R^n.


\end{multline} 


Здесь, как и всюду ниже, 


\begin{gather*} 


((\partial/\partial x)^2+h)^\beta\equiv\prod^n_{\nu=1} 


((\partial/\partial x_\nu)^2+h_\nu(x_\nu))^{\beta_\nu},\\


\gamma(\alpha;\beta)=(\sgn(\alpha_1-\beta_1),\dotsc, 


\sgn(\alpha_n-\beta_n));\ \; \sgn 0=0; 


\ \; \sgn t=1 \ \; \forall t>0. 


\end{gather*}  


Подчеркнем, что уравнения (7), в которых наибольшие мультииндексы  


имеют только четные координаты, всегда представимы в виде (5.2). 


 


Удобно обозначить 


$(\Lambda_k)^\alpha\equiv\prod\limits^n_{\nu=1} 


(\lambda_{k,\nu})^{\alpha_\nu}$, $k=(k_1,\dotsc,k_n)$. 


 


\begin{lems} 


При каждом $\varepsilon>0$, некоторой $c=c(\varepsilon)$ и всех 


$v(x)\in H$, 


удовлетворяющих условиям 


$\partial/\partial x_\nu v(x)\in H$, $|v(x)|\to 0$ при $|x|\to+\infty$, 


выполняются оценки 


\begin{equation} 


\|\partial/\partial x_\nu v(x)\|\le\varepsilon\| 


((\partial/\partial x_\nu)^2+h_\nu(x_\nu))v(x)\|+c\|v(x)\|. 


\end{equation} 


\end{lems} 


 


\begin{pf} 


Так как $\re h_\nu(t)\le\alpha_0$, то 


\begin{multline*} 


\|((\partial/\partial x_\nu)^2+h_\nu(x_\nu)-\alpha_0- 


\ov\varepsilon^{\,2})v(x)\|\,\|v(x)\|\ge \\ 


\ge|\la((\partial/\partial x_\nu)^2-\ov\varepsilon^{\,2}) 


v(x),\;v(x)\ra+\la(h_\nu(x_\nu)-\alpha_0)v(x),\;v(x)\ra|\ge \\ 


\ge |\la\partial/\partial x_\nu v(x),\ \partial/\partial x_\nu 


v(x)\ra+\ov\varepsilon^{\,2}\la v(x),\;v(x)\ra+ 


\la-(\re h_\nu(x_\nu)-\alpha_0)v(x),\;v(x)\ra|\ge \\ 


\ge \|\partial/\partial x_\nu v(x)\|^2+\ov\varepsilon^{\,2} 


\|v(x)\|^2. 


\end{multline*} 


Поэтому $\ov\varepsilon^{\,2}\|v(x)\|\le\|((\partial/\partial 


x_\nu)^2+h_\nu(x_\nu)-\alpha_0-\ov\varepsilon^{\,2})v(x)\|$. 


Учитывая последние 


две серии оценок, получим 


$\ov\varepsilon^{\,2}\|\partial/\partial x_\nu v(x)\|^2\le 


\|((\partial/\partial x_\nu)^2+h_\nu(x_\nu)-\alpha_0- 


\ov\varepsilon^{\,2})v(x)\|^2$. 


Отсюда легко следует, что 


$\varepsilon\|((\partial/\partial x_\nu)^2+h_\nu(x_\nu))v(x)\|+ 


\sqrt{|\alpha_0\varepsilon^2+1|} 


\|v(x)\|\ge\|\partial/\partial x_\nu v(x)\|$. 


\end{pf} 


 


\begin{lems} 


Пусть при всех $v(x)\in H$, удовлетворяющих условиям


$\partial/\partial x_\nu v(x)\in H$ $\forall\nu$, $|v(x)|\to 0$


при $|x|\to +\infty$, для 


каждого $\varepsilon>0$ выполняются  


оценки $(5.3)$, элементы матриц $\wt c_{\alpha,\beta}(x)$ измеримы и 


ограничены в существенном, 


для всех $\beta<\alpha$, $\alpha\in\Phi(2)$ 


\begin{equation} 


|(\Lambda_k)^\beta|\bigg(1+\sum_{\alpha\in\Phi(2)}\; 


\sum_{\gamma\le\alpha}|(\Lambda_k)^\gamma|\bigg)^{-1}\to 0 


\quad 


\text{если}\quad |k|\to+\infty. 


\end{equation} 


Тогда при каждом $\varepsilon_1>0$, некоторой $c=c(\varepsilon_1)$ и 


всех $\psi(x)\in L_\varphi$ 


\begin{multline} 


S_1\equiv\bigg\|\sum_{\alpha\in\Phi(2)}\; 


\sum_{\beta<\alpha}\wt c_{\alpha,\beta}(x)


(\partial/\partial x)^{\gamma(\alpha,\beta)} 


((\partial/\partial x)^2+h)^\beta\psi(x)\bigg\|\le \\ 


\le \varepsilon_1\bigg(\,\sum_{\alpha\in\Phi(2)}\; 


\sum_{\beta\le\alpha}\|((\partial/\partial x)^2+h)^\beta 


\psi(x)\|\bigg)+c\|\psi(x)\|. 


\end{multline} 


\end{lems} 


 


\begin{pf} 


Левая часть доказываемой оценки не больше 


$$ 


S_2\equiv B\sum_{\alpha\in\Phi(2)}\; 


\sum_{\beta<\alpha}\|(\partial/\partial x)^{\gamma(\alpha,\beta)} 


((\partial/\partial x)^2+h)^\beta\psi(x)\|, 


$$ 


причем $B$ не зависит от $\psi(x)$. Так как $\psi(x)\in L_\varphi$, то 


\begin{equation} 


((\partial/\partial x)^2+h)^\beta\psi(x)=\sum_k\psi(k) 


(\Lambda_k)^\beta\bigg(\prod^n_{\nu=1}y_{(k_\nu,\nu)} 


(x_\nu)\bigg), 


\end{equation} 


и ненулевых слагаемых конечное число. Поэтому 


$$ 


(\partial/\partial x)^{\gamma(\alpha,\beta)} 


((\partial/\partial x)^2+h)^\beta\psi(x)\in H\ \; \forall 


\beta\le\alpha\in\Phi(2), \ \; |((\partial/\partial x)^2 


+h)^\beta\psi(x)|\to 0\ \; 


\text{при}\ \; |x|\to+\infty, 


$$ 


ибо имеет место (5.1). 


Отсюда и из (5.3) следует 


$$ 


S_2\le\varepsilon\sum_{\alpha\in\Phi(2)}\; 


\sum_{\beta<\alpha}\|((\partial/\partial x)^2+ 


h)^{\gamma(\alpha,\beta)+\beta}\psi(x)\|+ 


c\sum_{\alpha\in\Phi(2)}\bigg(\, 


\sum_{\beta<\alpha}\|((\partial/\partial x)^2+h)^\beta 


\psi(x)\|\bigg). 


$$ 


С другой стороны, в силу (5.4), (5.6) имеем 


$$ 


c\sum_{\alpha\in\Phi(2)}\,


\sum_{\beta<\alpha}\|((\partial/\partial x)^2+h)^\beta\psi(x)\|\le


\varepsilon_2\sum_{\alpha\in\Phi(2)}\,


\sum_{\beta\le\alpha}\|((\partial/\partial x)^2+h)^\beta\psi(x)\|+ 


c(\varepsilon_2)\|\psi(x)\|,\ \;


0<\varepsilon_2<\varepsilon. 


$$ 


Полагая $\varepsilon_1=\varepsilon-\varepsilon_2$ и учитывая 


$\gamma(\alpha,\beta)+\beta\le\alpha$, получим (5.5). 


\end{pf} 


 


\begin{lems} 


Предположим, что $a_\alpha(x)\in C^{\{0\}}_{m,m}(\Bbb R^n)$, для любого 


$\varepsilon_2>0$, 


некоторого $\rho=\rho(\varepsilon_2)$ и всех $x$, $y$, удовлетворяющих 


условиям $\|x\|\ge\rho$, $\|y\|\ge\rho$, 


выполняются оценки 


\begin{equation} 


\sum_{\alpha\in\Phi(2)}|a_\alpha(x)-a_\alpha(y)|\le\varepsilon_2.


\end{equation} 


Тогда существуют такие скалярные вещественные функции $\omega_r(x)\in 


C^{+\infty}_1(\Bbb R^n)$, что 


$(\omega_1(x))^2+\dotsb+(\omega_M(x))^2\equiv1$, $\omega_M(\tau)=1$ 


при $|\tau|\ge\rho_2;$ для всех $x\in\supp(\omega_r(x))$, 


$y\in\supp(\omega_r(x))$ и 


любого $r$ выполняются оценки $(5.7)$. 


\end{lems} 


 


\begin{pf} 


Множества $\Omega_r=\{x:|x-\tau(r)|<\rho_1\}$ выберем  


так, чтобы $\{|x|\le\rho\}\subset\Omega_1\cup\dotsb\cup\Omega_{M-1}$ и 


при всех $r$, $x\in\ov\Omega_r$, $y\in\ov\Omega_r$ 


выполнились оценки (5.7). При $r=1,\dotsc,M-1$ положим 


$\wt f_r(x)=\exp(|x-\tau(r)|^2-p^2_1)$, 


если $x\in\Omega_r$ и $\wt f_r(x)=0$ для $x\notin\Omega_r$;


$\wt f_M(x)=\exp((\rho^2-|x|^2)^{-1})$ при $|x|>\rho$ и 


$\wt f_M(x)=0$ при $|x|\le\rho$. 


Так как $\{|x|\le\rho\}\subset\Omega_1\cup\dotsb\cup\Omega_{M-1}$, то 


$s(x)=((\wt f_1(x))^2+\dotsb+(\wt f_{M-1}(x))^2+ 


(\wt f_M(x))^2)^{1/2}>0$. 


Поэтому функции  


$\omega_r(x)=\wt f_r(x)/s(x)$ принадлежат $C^{+\infty}_1(\Bbb R^n)$ и 


$(\omega_1(x))^2+\dotsb+(\omega_M(x))^2=1$. 


Если $x\in\supp\omega_r(x)$, $y\in\supp\omega_r(x)$ и $r\le M-1$, то 


$x\in\Omega_r$, $y\in\Omega_r$, и выполняется (5.7). 


В случае $x,y\in\supp\omega_M(x)$ имеем $|x|\ge\rho$, $|y|\ge\rho$, 


поэтому  


выполняются (5.7). Фиксируем $\rho_2$ так, чтобы 


$\{|x|\ge\rho_2\}\supset\Omega_1\cup\dotsb\cup\Omega_{M-1}$. 


Для $|x|\ge\rho_2$ имеем 


$x\notin\Omega_1\cup\dotsb\cup\Omega_{M-1}$, $\omega_r(x)=0$ 


$\forall r\le M-1$, т.\,е. 


$\omega_M(x)=1$. 


\end{pf} 


 


Удобно через $C_\nu((\cal D^2+h)^\beta)$ обозначить множество всех тех 


функций $a(x)\in C^{\{\alpha\}}_\nu(\Bbb R^n)$, 


для которых при всех $w(x)\in C^{+\infty}_1(\Bbb R^n)$ в равенствах 


$$ 


((\partial/\partial x)^2+h)^\beta(w(x)a(x))\equiv 


\sum_{\delta\le\beta}C_{\delta,\beta}(x) 


(\partial/\partial x)^{\gamma(\delta,\beta)} 


((\partial/\partial x)^2+h)^\delta w(x) 


$$ 


функции $C_{\delta,\beta}(x)$, не зависящие от $w(x)$, ограничены. 


Условие $a(x)\in C_\nu((\cal D^2+h)^\beta)$ 


означает, что при $\gamma\ne0$ производные $a^{(\gamma)}(x)$ достаточно 


быстро 


убывают (естественно при условии $\gamma\le\beta$). 


 


\begin{thms} 


Предположим, что при всех $\alpha\in\Phi(2)$, $\beta\in\Phi(2)$ 


выполняются 


включения $a_\alpha(x)\in C_{m,m}((\cal D^2+h)^{2\beta})$ и функции 


$a_\alpha(x)$ имеют пределы при  


$|x|\to+\infty;$ элементы матриц $B_{\alpha,\beta}(x)$ измеримы и 


ограничены в существенном$;$ 


при некоторых $g_1>0$, $t\in\Bbb C$ и всех $x\in\Bbb R^n$, $k\in\Bbb 


Z^n$, $z\in\Bbb C^n$ имеют место  


оценки 


\begin{equation} 


\bigg|\bigg(t E+\sum_{\alpha\in\Phi(2)}a_\alpha(x) 


(\Lambda_k)^\alpha\bigg)z\bigg|^2\ge g_1\bigg[(d(t))^2+ 


\sum_{\alpha\in\Phi(2)}\;


\sum_{\beta\le\alpha}|(\Lambda_k)^\beta|^2\bigg]|z|^2, 


\end{equation} 


в которых $d(t)\to+\infty$ при $|t|\to+\infty;$ для всех 


$\beta<\alpha$, $\alpha\in\Phi(2)$ и $|k|\to+\infty$ 


$$ 


\bigg(|(\Lambda_k)^\beta|\bigg(1+ 


\sum_{\alpha\in\Phi(2)}\;


\sum_{\gamma\le\alpha}|(\Lambda_k)^\gamma|\bigg)^{-1}\bigg)\to 0; 


$$ 


при каждом $\varepsilon>0$ и всех функциях $v(x)\in H$, удовлетворяющих 


условиям 


$(\partial/\partial x_\nu v(x))\in H$ $\forall\nu;$ $|v(x)|\to0$ при 


$|x|\to+\infty$, выполняются оценки $(5.3)$. 


 


Тогда 


\begin{itemize} 


\item[1.] спектр $\varphi$-задачи для уравнения $(7)$ точечен$;$ 


\item[2.] если $\varphi$-решение $u$ уравнения $(7)$ существует, то оно


удовлетворяет условию 


\end{itemize} 


$$ 


\sum_k 


\sum_{\alpha\in\Phi(2)}\;


\sum_{\beta\le\alpha} 


|(\Lambda_k)^\beta|^2|u_k|^2<+\infty. 


$$ 


\end{thms} 


 


\begin{pf} 


На первом этапе докажем теорему для частного 


случая, когда $B_{\alpha,\beta}(x)$ --- нулевые матрицы. Через $\cal 


L_t(y)$ обозначим 


дифференциальный полином 


$$ 


t E+\sum_{\alpha\in\Phi(2)}a_\alpha(y)((\partial/\partial x)^2 


+h)^\alpha; 


$$ 


через $q(z)$ --- скалярный полином, удовлетворяющий условию 


$$ 


|q(\Lambda_k)|\ge 


\sum_{\alpha\in\Phi(2)}\;\sum_{\beta\le\alpha} 


|(\Lambda_k)^\beta|\quad \forall k\in\Bbb Z^n; 


$$ 


через $Q$ --- $\varphi$-стационарный оператор, определяемый равенствами 


$$ 


Q w=\sum_k q(\Lambda_k)w_k\bigg(\prod^n_{j=1} 


y_{(k_j,j)}(x_j)\bigg). 


$$ 


Отметим, что в силу (5.8) за дифференциальный полином 


$q(\partial/\partial x)$ можно  


взять полином $\cal L_t(y)e(1)/g_1$, в котором $|t|$ достаточно велико. 


 


Фиксируем малое число $\varepsilon>0$ и функции 


$\omega_1(x),\dotsc,\omega_M(x)$, 


построенные в лемме 5.3. Тогда для $\psi(x)\in L_\varphi$ имеем 


$\|\cal L_t(y)\psi(x)\|\ge F_1-F_2-F_3$, 


где 


\begin{align*} 


F_1&=\bigg(\sum^M_{\nu=1}\|\cal L_t(\tau(\nu))(\omega_\nu(x)\psi(x) 


)\|^2\bigg)^{1/2};\\ 


F_2&=\bigg(\sum^M_{\nu=1}\|(\cal L_0(x)-\cal L_0(\tau(\nu))) 


(\omega_\nu(x)\psi(x))\|^2\bigg)^{1/2};\\ 


F_3&=\bigg(\sum^M_{\nu=1}\|\omega_\nu(x)\cal L_0(x)\psi(x) 


-\cal L_0(x)(\omega_\nu(x)\psi(x))\|^2\bigg)^{1/2}. 


\end{align*} 


Так как имеет место (5.7), $\tau(\nu)\in\Omega_\nu$ и 


$$ 


\|(\cal L_0(x)-\cal L_0(\tau(\nu))) 


(\omega_\nu(x)\psi(x))\|^2=\int_{\Omega_\nu} 


|(\cal L_0(x)-\cal L_0(\tau(\nu))) 


(\omega_\nu(x)\psi(x))|^2dx, 


$$ 


то 


\begin{multline*} 


F_2\le\varepsilon\bigg(\sum^M_{\nu=1} 


\sum_{\alpha\in\Phi(2)}\|((\partial/\partial x)^2+ 


h)^\alpha(\omega_\nu(x)\psi(x))\|^2\bigg)^{1/2}\le \\ 


\le\varepsilon\bigg(\sum^M_{\nu=1} 


\sum_{\alpha\in\Phi(2)}\|\omega_\nu(x) 


((\partial/\partial x)^2+h)^\alpha\psi(x)\|^2\bigg)^{1/2}+ \\ 


+\varepsilon\bigg(\sum^M_{\nu=1} 


\sum_{\alpha\in\Phi(2)}\|((\partial/\partial x)^2+h)^\alpha 


(\omega_\nu(x)\psi(x))-\omega_\nu(x) 


((\partial/\partial x)^2+h)^\alpha\psi(x)\|^2\bigg)^{1/2}. 


\end{multline*} 


В силу выбора $\omega_1(x),\dotsc,\omega_M(x)$ первая группа слагаемых 


совпадает с 


$$ 


F_{2,1}\equiv\varepsilon\bigg(\sum_{\alpha\in\Phi(2)} 


\|((\partial/ x)^2+h)^\alpha\psi(x)\|^2\bigg)^{1/2}. 


$$ 


К оценке второй группы слагаемых применима лемма 5.2. Поэтому 


\begin{multline} 


F_2\le\varepsilon\bigg( 


\sum_{\alpha\in\Phi(2)}\|((\partial/\partial x)^2+h)^\alpha 


\psi(x)\|^2\bigg)^{1/2}+ \\


+\varepsilon_1\bigg( 


\sum_{\alpha\in\Phi(2)}\;\sum_{\beta\le\alpha} 


\|((\partial/\partial x)^2+h)^\beta\psi(x)\|^2\bigg)^{1/2} 


+c\|\psi(x)\|. 


\end{multline} 


К оценке $F_3$ также можно применить лемму 5.2, т.\,е. 


$$ 


F_3\le\varepsilon\bigg( 


\sum_{\alpha\in\Phi(2)}\;\sum_{\beta\le\alpha}\|(( 


\partial/\partial x)^2+h)^\beta\psi(x)\|^2\bigg)^{1/2}+c\|\psi(x)\|. 


$$ 


Таким образом, доказано, что 


\begin{equation} 


F_2+F_3\le 3\varepsilon\bigg( 


\sum_{\alpha\in\Phi(2)}\;\sum_{\beta\le\alpha} 


\|((\partial/\partial x)^2+h)^\beta\psi(x)\|^2\bigg)^{1/2}+ 


c\|\psi(x)\|. 


\end{equation} 


Напомним, что несущественные постоянные мы обозначаем одними и теми


же буквами. Труднее оценить снизу величину $F_1$. Так как 


$(\omega_\nu(x)\psi(x))^{(\beta)}\to 0$ 


для всех $\beta$ и $\nu$, если $|x|\to+\infty$, то 


$$ 


\la(\omega_\nu(x)\psi(x))^{(\beta)},\ e(j)y_{(k)}(x)\ra= 


(-1)^{[\beta]}\la\omega_\nu(x)\psi(x),\ e(j) 


y^{(\beta)}_{(k)}(x)\ra. 


$$ 


Поэтому 


\begin{multline*} 


\la((\partial/\partial x)^2+h)^\beta(\omega_\nu(x)\psi(x)), 


\ e(j)y_{(k)}(x)\ra= \\ 


=\la\omega_\nu(x)\psi(x),\ ((\partial/\partial x)^2 


+\ov h)^\beta(e(j)y_{(k)}(x))\ra=(\Lambda_k)^\beta 


\la\omega_\nu(x)\psi(x),\ e(j)y_{(k)}(x)\ra. 


\end{multline*} 


Из выписанных равенств и ортонормированности $\varphi$, во-первых, 


следует сходимость по норме рядов 


$$ 


\sum_k(\Lambda_k)^\beta z_k y_{(k)}(x)\quad 


\forall\beta\le\alpha,\ \; \forall\alpha\in\Phi(2), 


$$ 


в которых $z_k$ --- вектор-столбец с координатами 


$z_{k,j}=\la\omega_\nu(x)\psi(x),\ e(j)y_{(k)}(x)\ra$, $j=1,\dotsc,m$, 


во-вторых, справедливо равенство 


$$ 


F_1=\sum^M_{\nu=1}\bigg(\sum_k\bigg|\bigg(t E+ 


\sum_{\alpha\in\Phi(2)}a_\alpha(\tau(\nu))(\Lambda_k)^\alpha 


\bigg)z_k\bigg|^2\bigg)^{1/2}. 


$$ 


Учитывая (5.8), получим 


\begin{equation} 


F_1\ge\bigg(g_1\sum^M_{\nu=1} 


\sum_{\alpha\in\Phi(2)}\;\sum_{\beta\le\alpha} 


\|((\partial/\partial x)^2+h)^\beta 


(\omega_\nu(x)\psi(x))\|^2+ 


g_1(d(t))^2\|\omega_\nu(x)\psi(x)\|^2\bigg)^{1/2}. 


\end{equation} 


Так как 


\begin{multline*} 


\|((\partial/\partial x)^2+h)^\beta(\omega_\nu(x)\psi(x))\|^2 


\ge \frac{1}{2}\|\omega_\nu(x)((\partial/\partial x)^2 


+h)^\beta\psi(x)\|^2 - \\ 


-\|\omega_\nu(x)((\partial/\partial x)^2+h)^\beta\psi(x)- 


((\partial/\partial x)^2+h)^\beta 


(\omega_\nu(x)\psi(x))\|^2, 


\end{multline*} 


то в силу леммы 5.2 при каждом $\varepsilon_1>0$, не зависящим от 


$\varepsilon$, 


\begin{multline*} 


\|((\partial/\partial x)^2+h)^\beta 


(\omega_\nu(x)\psi(x))\|^2\ge\frac{1}{2} 


\|\omega_\nu(x)((\partial/\partial x)^2+h)^\beta\psi(x)\|^2- \\ 


-\varepsilon_1\sum_{\gamma\le\beta} 


\|((\partial/\partial x)^2+h)^\gamma\psi(x)\|^2- 


c\|\psi(x)\|^2. 


\end{multline*} 


Отсюда и из (5.11) следует неравенство 


$$ 


F^2_1\ge\big(\tfrac{g_1}{2}-\varepsilon_1B)\bigg( 


\sum_{\alpha\in\Phi(2)}\;\sum_{\beta\le\alpha} 


\|((\partial/\partial x)^2+h)^\beta\psi(x)\|^2\bigg)+ 


((d(t))^2-c)\|\psi(x)\|^2. 


$$ 


В нем $B$ не зависит от $\varepsilon_1$, но зависит от $\varepsilon$. 


Таким образом, 


доказано, что в случае $(d(t))^2\ge c$ имеет место 


$$ 


F_1-(F_2+F_3)\ge(g_1/2-\varepsilon_1B-3\varepsilon)\bigg( 


\sum_{\alpha\in\Phi(2)}\;\sum_{\beta\le\alpha} 


\|((\partial/\partial x)^2+h)^\beta\psi(x)\|^2\bigg)^{1/2}+ 


((d(t))^2-c)^{1/2}\|\psi(x)\|. 


$$ 


Здесь $\varepsilon$, $\varepsilon_1$ --- положительные числа, которые 


можно выбрать сколь угодно 


малыми; $g_1$, $B$, $c$ не зависят от $\psi(x)$ и от $t$; $g_1$ не 


зависит от $\varepsilon$, $\varepsilon_1$; 


$g_1$, $B$ не зависит от $\varepsilon_1$. С другой стороны, 


$$ 


\|Q\psi(x)\|^2=\sum_k 


\sum_{\alpha\in\Phi(2)}\;\sum_{\beta\le\alpha} 


|(\Lambda_k)^\beta|^2|\psi_k|^2. 


$$ 


Поэтому выполняются оценки 


\begin{equation} 


\|\cal L_t(x)\psi(x)\|\ge g_2\|Q\psi(x)\|+ 


((d(t))^2-c)^{1/2}\|\psi(x)\|. 


\end{equation} 


Следовательно, выполняются аналоги оценок (2.1). Обозначим через 


$\cal L^+_{\ov t}(x)$ 


оператор, формально сопряженный к $\cal L_t(x)$. В силу предположений 


относительно $a_\alpha(x)$ оператор $\cal L^+_{\ov t}(x)$ определяется 


равенствами 


\begin{multline*} 


\cal L^+_{\ov t}(x)\psi(x)= 


\sum_{\alpha\in\Phi(2)}(a_\alpha(x))^*((\partial/\partial x)^2+ 


\ov h)^\alpha\psi(x)+\ov t\psi(x)+ \\ 


+ \sum_{\alpha\in\Phi(2)}\;\sum_{\beta<\alpha} 


\wt B_{\alpha,\beta}(x)(\partial/\partial x)^{\gamma(\alpha,\beta)} 


((\partial/\partial x)^2+\ov h)^\beta\psi(x)\quad 


\forall\psi(x)\in L_\varphi, 


\end{multline*} 


причем элементы матриц $\wt B_{\alpha,\beta}(x)$ ограничены и измеримы. 


Так как 


матрицы $a_\alpha(x)$ конечномерны и квадратны, то в силу (5.8) 


\begin{equation} 


\bigg|\bigg(\ov t E+ 


\sum_{\alpha\in\Phi(2)}(a_\alpha(x))^*(\ov\Lambda_k)^\alpha\bigg) 


z\bigg|^2\ge g_1\bigg((d(t))^2+ 


\sum_{\alpha\in\Phi(2)}\;\sum_{\beta\le\alpha} 


|(\Lambda_k)^\beta|^2\bigg)|z|^2, 


\end{equation} 


причем $g_1$ и $d(t)$ --- объекты, используемые в (5.8). Оценки (5.13), 


как доказано ранее, гарантируют 


$$ 


\bigg\|\ov t E+ 


\sum_{\alpha\in\Phi(2)}(a_\alpha(x))^*((\partial/\partial x)^2+ 


\ov h)^\alpha\psi(x)\bigg\|\ge g_2\|Q\psi(x)\|- 


((d(t))^2-c)^{1/2}\|\psi(x)\|. 


$$ 


Отсюда, а также из леммы 5.2 следует справедливость аналогов оценок 


(2.8). Так как $\cal D_Q\subset H$, то аналоги включений (2.11) 


выполняются автоматически. 


 


Убедимся в том, что при каждом достаточно малом $c_3>0$ выполняются 


аналоги оценок (2.13). Пусть $\psi(x)\in L_\varphi$, 


$v(x)\in L_\varphi$. Тогда 


$$ 


S_3\equiv\la(Q^+)^{-1}\psi(x),\ (Q\cal L_t(x)Q^{-1}- 


\cal L_t(x))v(x)\ra=\la((Q^+)^{-1}\cal L^+_{\ov t}(x)- 


\cal L^+_{\ov t}(x)(Q^+)^{-1})\psi(x),\ v(x)\ra. 


$$ 


Подчеркнем, что возможность интегрирования по частям и исчезновение  


внеинтегральных членов при интегрировании по частям гарантируется 


предположениями относительно $a_\alpha(x)$, $y_{(k,\nu)}(x_\nu)$. 


Запишем полином $q(\xi)$ 


в виде $\sum\limits_{\alpha\in\Phi(2)}l_\alpha(\xi)^\alpha$. Очевидно, 


\begin{multline*} 


Q\cal L_t(x)Q^{-1}v(x)= 


\sum_{\Tilde\alpha\in\Phi(2)} 


\sum_{\alpha\in\Phi(2)} 


l_{\Tilde\alpha}((\partial/\partial x)^2+h)^{\Tilde\alpha} 


[a_\alpha(x)((\partial/\partial x)^2+h)^\alpha 


(Q^{-1}v(x))]= \\ 


%


= \sum_{\Tilde\alpha\in\Phi(2)} 


\sum_{\alpha\in\Phi(2)}l_{\Tilde\alpha} 


\sum_{\beta\le\Tilde\alpha}c_{\beta,\alpha,\Tilde\alpha}(x) 


(\partial/\partial x)^{\gamma(\Tilde\alpha,\beta)} 


((\partial/\partial x)^2+h)^{\beta+\alpha} 


Q^{-1}v(x). 


\end{multline*} 


Обозначение $\gamma(\wt\alpha,\beta)$ 


введено в начале п.\,5. В последних 


равенствах $a_\alpha(x)=c_{\beta,\alpha,\Tilde\alpha}(x)$, 


если $\wt\alpha=\beta$. Другими словами, сумма членов, для которых 


$\beta=\wt\alpha$, 


совпадает с $\cal L_t(x)v(x)$. Поэтому 


$$ 


S_3=\sum_{\Tilde\alpha\in\Phi(2)} 


\sum_{\alpha\in\Phi(2)}\ov l_{\Tilde\alpha} 


\sum_{\beta<\Tilde\alpha}\la(c_{\beta,\alpha,\Tilde\alpha} 


(x))^*(Q^+)^{-1}\psi(x),\ (\partial/\partial x 


)^{\gamma(\Tilde\alpha,\beta)}((\partial/\partial x)^2 


+h)^{\beta+\alpha}(Q^{-1}v(x))\ra. 


$$ 


Предположения относительно $a_\alpha(x)$, $y_{(k,\nu)}(t)$ 


позволяют интегрировать 


по частям $(2\beta)$ раз. При этом внеинтегральные члены исчезнут 


и получим 


$$ 


S_3=\sum_{\Tilde\alpha\in\Phi(2)} 


\sum_{\alpha\in\Phi(2)}\ov l_{\Tilde\alpha} 


\sum_{\beta<\Tilde\alpha}\sum_{\delta\le\beta}\la 


c_{\beta,\delta,\alpha,\Tilde\alpha}(x) 


(\partial/\partial x)^{\gamma(\delta,\beta)} 


((\partial/\partial x)^2+\ov h)^\delta(Q^+)^{-1}\psi(x), 


\ ((\partial/\partial x)^2+h)^\alpha(Q^{-1}v(x))\ra, 


$$ 


причем требования $a_\alpha(x)\in C_{m,m}((\cal D^2+h)^\beta)$ 


гарантируют ограниченность элементов 


матриц $c_{\beta,\delta,\alpha,\Tilde\alpha}(x)$. Поэтому 


$$ 


|S_3|\le B\sum_{\Tilde\alpha\in\Phi(2)} 


\sum_{\alpha\in\Phi(2)} 


\sum_{\beta<\Tilde\alpha}\sum_{\delta\le\beta} 


\|(\partial/\partial x)^{\gamma(\delta,\beta)}((\partial/ 


\partial x)^2+\ov h)^\delta(Q^+)^{-1}\psi(x)\|\,\| 


((\partial/\partial x)^2+h)^\alpha(Q^{-1}v(x))\|. 


$$ 


В силу (5.5), (5,6) имеем 


\begin{align*} 


\|(\partial/\partial x)^{\gamma(\delta,\beta)}((\partial/ 


\partial x)^2+\ov h)^\delta(Q^+)^{-1}\psi(x)\|&\le 


B_1\|((\partial/\partial x)^2+h)^{\gamma(\delta,\beta)+\beta} 


(Q^+)^{-1}\psi(x)\|;\\ 


%


\|((\partial/\partial x)^2+h)^\alpha(Q^{-1}v(x))\|&\le  


B_2\|v(x)\|. 


\end{align*} 


Из последних трех серий оценок следует 


\begin{equation} 


|S_3|\le\wt B\sum_{\Tilde\alpha\in\Phi(2)}\;\sum_{\beta<\Tilde\alpha} 


\|((\partial/\partial x)^2+\ov h)^\beta(Q^+)^{-1}\psi(x)\|\, 


\|v(x)\|. 


\end{equation} 


При этом постоянная $\wt B$ не зависит от $\psi(x)$ и от $v(x)$. Если 


$\beta<\wt\alpha$, то в 


силу (5.4), (5.5) имеем 


\begin{multline*} 


\|((\partial/\partial x)^2+\ov h)^\beta(Q^+)^{-1}\psi(x)= 


\bigg(\sum_k|\psi_k|^2|(\Lambda_k)^\beta|^2 


|q(\Lambda_k)|^{-2}\bigg)^{1/2}\le \\ 


\le \varepsilon_1\bigg(\sum_k|\psi_k|^2\bigg)^{1/2}+ 


c\bigg(\sum_{|k|\le N}|\psi_k|^2\bigg)^{1/2}\le 


\varepsilon_1\|\psi(x)\|+c\|\psi(N,x)\|. 


\end{multline*} 


Итак, доказано, что 


$$ 


|S_3|\le(\wt B\varepsilon_1\|\psi(x)\|+\wt c\|\psi(N,x)\|) 


\|v(x)\|, 


$$ 


причем $\wt B$ не зависит от $\varepsilon_1$, и $\varepsilon_1$ можно 


выбрать сколь угодно малым.  


Так как в последней оценке постоянные не зависят от $v(x)$, и 


множество $L_\varphi$ плотно в $H$, то при $c_3=\wt B\varepsilon_1$ 


выполняется аналог оценок (2.13). 


В силу (5.12), (5.13) при больших $|t|$ числа $t$, $\ov t$ являются 


несобственными 


значениями соответственно операторов $\cal L_t(x)$, 


$\cal L^+_{\ov t}(x)$. При этом  


выполняются аналоги предположений теоремы 2.2. Поэтому последние  


числа $t$, $\ov t$ являются регулярными значениями соответственно 


операторов 


$\cal L_t(x)$, $\cal L^+_{\ov t}(x)$. Отсюда, из (5.12) и из леммы 5.2 


легко следует, что каждое несобственное значение $\varphi$-задачи для 


уравнения (7) является 


регулярным. Первое утверждение теоремы доказано. Второе --- 


простое следствие оценок (5.8). 


\end{pf} 


 


\begin{notes} 


Среди предположений теоремы 5.1 наиболее 


надуманным является предположение 


$a_\alpha(x)\in C_{m,m}((\cal D^2+h)^\beta)$. Естественно 


возникает 


вопрос: нельзя ли ограничиться непрерывностью $a_\alpha(x)$? В общем  


случае нельзя. Это объясняется тем, что в случае 


$\re h_\nu(t)=h_\nu(t)$ 


функции $h_\nu(t)$ обязаны быть неограниченными, что следует из теоремы 


А.М.\,Молчанова ([3], c.\,393). Поэтому отказ от последнего включения 


хотя бы при одном $\alpha$ 


привел бы к тому, что при исследовании спектра 


$\varphi$-задачи для уравнения (7) члены вида 


$b_{\alpha,\beta}(x)=(\partial/\partial x)^{\gamma(\alpha,\beta)} 


((\partial/\partial x)^2+h)^\beta$ при $\beta<\alpha$ 


играли бы существенную роль. 


\end{notes} 


 


\begin{notes} 


Мы не доказываем индивидуальную теорему о  


гладкости. Ее доказательство мало чем отличается от доказательства  


аналогичной теоремы для $(b{-}a)$-периодической задачи. Последняя  


будет доказана в третьей части. Для $\varphi$-задачи, рассматриваемой в  


части~II, индивидуальная теорема о гладкости гласит: 


\begin{itemize} 


\item[\phantom{m}] пусть выполняются предположения теоремы 5.1 и при 


$\gamma$, не превосходящем некоторого $\wt\beta\in\Phi(2)$, включения 


\begin{gather*} 


a_\alpha(x)\in C_{m,m}((\cal D^2+h)^{2(\beta+\gamma)});\ \; 


b_{\alpha,\beta}(x)\in C_{m,m}((\cal D^2+h)^\gamma)\ \; \forall 


\alpha\in\Phi(2),\ \;\forall\beta\in\Phi(2); \\


((\partial/\partial x)^2+h)^\gamma f(x)\in L^2_m(\Bbb R^n); 


\end{gather*}  


тогда в случае существования $\varphi$-решение $u$ уравнения (7) 


удовлетворяет условию 


$$ 


\sum_{\alpha\in\Phi(2)}\;\sum_{\beta\le\alpha}\sum_k 


|(\Lambda_k)^{\beta+\gamma}|^2|u_k|^2<+\infty. 


$$ 


\end{itemize} 


Напомним, что $u_k$ --- вектор-столбец из соответствующих коэффициентов  


Фурье $\varphi$-решения $u$. 


\end{notes} 


 


\begin{thebibliography}{9} 


 


\bibitem{1} 


Мокейчев В.С., \fbox{Мокейчев А.В.} {\em Новый подход к теории 


линейных задач для систем дифференциальных уравнений в частных 


производных}. I. // Изв. вузов. Математика. -- 1999. -- \No\,1. -- 


C.\,25--35. 


\bibitem{2} 


Мокейчев В.С. {\em Дифференциальные уравнения с отклоняющимися 


аргументами}. -- Казань: Изд-во Казанск. ун-та, 1985. -- 222 с. 


\bibitem{3} 


Наймарк М.А. {\em Линейные дифференциальные операторы}. -- 2-е изд.


-- М.: Наука, 1969. -- 526 с. 


 


\end{thebibliography} 


 


\vskip 2cm 


 


\post{\it Казанский государственный}{\it Поступила} 


\post{\it университет}{18.03.1996} 


 


\label{end} 


\end{document} 


